% Generated draft: requires mathematical and source review.
\documentclass[11pt]{article}
\usepackage[margin=25mm]{geometry}
\usepackage{fontspec}
\setmainfont{Latin Modern Roman}
\usepackage{amsmath,amssymb,mathtools,unicode-math}
\setmathfont{Latin Modern Math}
\usepackage{enumitem,needspace,xurl,hyperref}
\setlength{\parindent}{0pt}
\setlength{\parskip}{5pt}
\setlength{\emergencystretch}{3em}
\newcommand{\N}{\mathbb{N}}
\newcommand{\Z}{\mathbb{Z}}
\newcommand{\Q}{\mathbb{Q}}
\newcommand{\R}{\mathbb{R}}
\newcommand{\C}{\mathbb{C}}
\newcommand{\F}{\mathbb{F}}
\newcommand{\A}{\mathbb{A}}
\newcommand{\PP}{\mathbb{P}}
\newcommand{\E}{\mathbb{E}}
\newcommand{\eps}{\varepsilon}
\newcommand{\dd}{\,\mathrm{d}}
\begin{document}
\section*{UnsolvedMath ID: 30000234}
\textbf{Source problem code:} catalogue entry
\subsection{Definitions and mathematical statement}
Let \(\F_q\) be a finite field, where \(q\) is a prime power, and let \(A\subseteq\F_q\) be nonempty. Define
\[
 A+A=\{a+b:a,b\in A\},\qquad
 A\cdot A=\{ab:a,b\in A\}.
\]
The source asks whether an Erd\H{o}s--Szemer\'edi sum-product theorem holds in finite fields under an appropriate exclusion of nontrivial subfields. It does not specify the quantitative bound, allowable sizes of \(A\), or the exclusion condition.

A customary superlinear-growth form would seek constants \(C>0\) and \(\delta>0\) for which
\[
 \max\{|A+A|,|A\cdot A|\}\geq C|A|^{1+\delta}
\]
for every pair \((\F_q,A)\) satisfying precise size and subfield-nonconcentration hypotheses. This is an illustrative template only: the constants, their dependencies, the exponent and the admissible class must be recovered from the report. The extracted name alone does not identify this fixed-exponent target with a nearly quadratic bound for every exponent below two, or with a particular error term or endpoint range.

Subfields supply a direct obstruction: for a subfield \(K\subseteq\F_q\), one has \(K+K=K\) and \(K\cdot K=K\). Thus a restriction is essential, but ``exclusion of nontrivial subfields'' does not specify whether it excludes only \(A\) equal to or contained in a proper subfield, controls large intersections with subfields, or also concerns translated or multiplicatively scaled copies. The meaning of nontrivial, including treatment of the prime subfield, is unstated. These alternatives are not silently identified.

There is a separate ambient-size issue: both displayed sets have size at most \(q\), so a bound \(C|A|^{1+\delta}\) with fixed positive \(C,\delta\) cannot hold for unrestricted \(A=\F_q\) as \(q\) grows, even for prime fields with no proper subfields. A suitable size range or a saturating bound must therefore be specified in addition to the intended subfield condition; no particular correction is inserted.

The literature triage checked 2026-08-21 says the report poses the finite-field analogue and that its checked literature gives no general affirmative or negative resolution. With the target and hypotheses unspecified, this is retained as dated curation metadata, not an independent claim that all finite-field sum-product estimates are open. The verification field reports source checking and normalized wording but supplies no missing quantitative formulation.
\subsection{Short English statement}
The question seeks finite-field sum-product growth once subfield obstructions are appropriately excluded. To state a definite theorem, one must specify the growth exponent, the size range of the set and a precise nonconcentration condition. The ambient field-size bound is an additional constraint independent of proper subfields.
\subsection{Sources}
\begin{itemize}
\item Local UnsolvedMath registry, permanent ID 30000234, source code catalogue entry, ``Finite-Field Erd\H{o}s--Szemer\'edi Sum-Product Bounds,'' including the original and cleaned question and verification metadata.
\item \emph{Complexity Theory (2006)}, Oberwolfach Reports, pp.\ 1431--1504, DOI 10.4171/owr/2005/26: \url{https://doi.org/10.4171/owr/2005/26}. This is the sole supplied literature reference, labelled ``OWR report.'' The citation year is 2006, while the DOI and proposed-year metadata use 2005.
\item Literature triage checked 2026-08-21; record updated 2026-09-16. No external source has been visited or independently reviewed here.
\end{itemize}
\end{document}
